% ----------------------
\subsection{Section 126 (9 July 1841)}
% ----------------------
	\label{DC126}
	
	% -----
	\subsubsection{Historical Background}
	% -----
		\label{DC126-HB}
		
		% --
		\paragraph{JSP}
		% --
		
			``On 9 July 1841, JS dictated a revelation for Brigham Young, releasing him from extended travel and admonishing him to remain with and care for his family. From the time Young joined the church in 1832, missionary travels marked his life. Young had departed Nauvoo, Illinois, almost two years earlier for a mission to Europe with other apostles and church members, leaving his family behind. He and the others proselytized primarily in England, where thousands converted to the church. While there, Young's fellow apostles sustained him as the president of their quorum. Young and the other apostles serving missions were deeply concerned about the well-being of their families. During Young's sojourn in England, his wife, Mary Ann Angell Young, and their children struggled with poverty, lacked essential goods, and suffered from illness. Despite her downtrodden circumstances, Mary Ann expressed gratitude for her husband's efforts. `I am glad to hear the work of the Lord is prospering in england,' she wrote to Young in April 1841. `It gives me much joy.'
			
			``Young and his fellow apostles Heber C. Kimball and John Taylor returned to Nauvoo, arriving on 1 July 1841. There they `were met by President Smith, and many . . . old tried friends, whose countenances expressed the most heartfelt satisfaction at [their] return.' Eight days later, JS dictated the revelation featured here at Young's residence in Nauvoo. The revelation acknowledged Young's missionary labors away from home and commanded him thereafter to `send' the gospel message abroad---implying that he should reside at home---and to provide `special care' for his family.
			
			``The revelation's direction that Young no longer leave his family for extended periods was likely also connected to Young's increasing responsibilities in Nauvoo. A month after dictating this revelation, JS declared in a public discourse that `the time had come when the twelve should be called upon to stand in their place next to the first presidency, and attend to the settling of emegrants and the business of the church at the stakes.' JS further stated that the Twelve Apostles had earned the right to stay with their families, where they would have better opportunity to provide for them.
			
			``The original copy of this revelation has not been located. It is likely that a loose dictation copy was created, after which a copy was given to Young and the original was kept in JS's office. Willard Richards copied the text of this revelation into the Book of the Law of the Lord on 17 December 1841, four days after he was appointed recorder for the temple and JS's personal scribe.''
					
	% -----
	\subsubsection{Versions}
	% -----
		\label{DC126-V}
		
		See the \href{https://drive.google.com/file/d/1goAvNz8jS4R8slYfMG_db55Ty2z_vmgK/view?usp=sharing}{sections comparison} document.
	
		\begin{compactdesc}
			\item[JSP] \hyperref[EMS-1841-JS-9-JUL-1841]{9 July 1841}
			\item[BC] ---
			\item[DC 1835] ---
			\item[DC 1844] ---
			\item[DN] 5:6 (18 April 1855), 41
			\item[MS] 34:11 (12 March 1872), 165
		\end{compactdesc}

	% -----
	\subsubsection{Section 126:1} % *DC126-1 
	% -----
		\label{DC126-1}

		DEAR and well-beloved brother, Brigham Young, verily thus saith the Lord unto you: My servant Brigham, it is no more required at your hand to leave your family as in times past, for your offering is acceptable to me.

		% \begin{framed}
		% \scriptsize
			% \begin{compactdesc}
				% \item[JSP] 
				% % \item[BC] 
				% % \item[]
			% \end{compactdesc}
		% \end{framed}

	% -----
	\subsubsection{Section 126:2} % *DC126-2 
	% -----
		\label{DC126-2}

		I have seen your labor and toil in journeyings for my name.

		% \begin{framed}
		% \scriptsize
			% \begin{compactdesc}
				% \item[JSP] 
				% % \item[BC] 
				% % \item[]
			% \end{compactdesc}
		% \end{framed}

	% -----
	\subsubsection{Section 126:3} % *DC126-3 
	% -----
		\label{DC126-3}

		I therefore command you to send my word abroad, and take especial care of your family from this time, henceforth and forever.  Amen.

		% \begin{framed}
		% \scriptsize
			% \begin{compactdesc}
				% \item[JSP] 
				% % \item[BC] 
				% % \item[]
			% \end{compactdesc}
		% \end{framed}